\scanpage{38}
Il y a lieu de \unsure{s'intéresser|travailler}
$1^\circ$) au cas où on a un anneau
de base $k$ de Dedekind,
et $2^\circ$) \unsure{au|un|avec} $k$ quelconque \inserted{(\unsure{évt comm})}
(donc $\unsure{\Theta|\mathcal C|C}\subset\widehat\Delta_{k\text{-}\mathrm{mod}}$). La \deleted{\unsure{distinction|situation}}
conclusion à retenir, c'est
qu'on connaît les Hom de
$X_*$ dans tt $k$-module str.
$\partial$-\deleted{\unsure{parfaits}}\inserted{parfait $M$ (comme foncteur de $M$)} dès qu'on connaît
le complexe \inserted{de cochaînes ``non dég''} $C^{\unsure{*!|*}}(X_*,k)$ à coeff. dans
$k$. Notre propos principal
est de voir comment reconstituer
la structure multiplicative
et ``puiss. divisées'') sur le
foncteur
\[
 M\mapsto F_{X_*}(M_*)=F(M_*)=\operatorname{Hom}_{\widehat\Delta}(X_*,M_*)
\]
i.e. les homs
\[
 \left\{\begin{aligned}
 F(M_*)\otimes_k F(N_*)&\longrightarrow F(M_*\otimes N_*)\\
 \Gamma^i F(\deleted{\unsure{\Gamma^i}}M_*)&\longrightarrow F(\Gamma^i(M^{\unsure{n|*|\bullet}}))
 \end{aligned}\right.
\]
en tant que \deleted{\illegible{1}} structures
sur le complexe \inserted{de cochaînes} $C^*=C^{*!}(X_*,k)$

Il faudrait examiner si pour
$k$ de car 0 \inserted{$C^*$ avec} cette structure suppl.
% archive sheet number: 37
